\documentclass[11pt]{article}
\usepackage[margin=1in]{geometry}
\usepackage{amsmath,amssymb,amsthm,mathtools}
\usepackage{graphicx}
\usepackage[hidelinks]{hyperref}

\newtheorem{theorem}{Theorem}
\newtheorem{lemma}{Lemma}
\newcommand{\C}{\mathbb C}
\newcommand{\M}{\mathbb M}
\newcommand{\OS}{\mathcal{OS}}
\newcommand{\K}{\mathcal K}
\newcommand{\G}{\mathcal G}
\newcommand{\Rea}{\operatorname{Re}}
\newcommand{\Ima}{\operatorname{Im}}
\newcommand{\mconv}{\operatorname{mconv}}
\newcommand{\norm}[1]{\left\lVert #1\right\rVert}

\title{A Contractive but Not Completely Positive Map on \(\OS(J_3)\)}
\author{\texttt{agent\_lane\_03}}
\date{2026-07-22}

\begin{document}
\maketitle

\begin{center}
\fbox{\parbox{0.91\linewidth}{\textbf{Status.}
Candidate counterexample, likely valid. The construction is existential but
finite-dimensional: the separating operator can be chosen in \(\M_3\).}}
\end{center}

\section*{Source and Conjecture}

The source is Pankaj Dey, Atanu Dhang, and Mithun Mukherjee,
\emph{Matricial ranges, dilations, and unital contractive maps},
arXiv:2509.18027 \cite{DDM}. Write
\[
 J_3=\begin{bmatrix}0&1&0\\0&0&1\\0&0&0\end{bmatrix},
 \qquad \OS(J_3)=\operatorname{span}\{I_3,J_3,J_3^*\}.
\]
Conjecture 6.10, on page 21, states that for a unital self-adjoint
preserving map \(\varphi:\OS(J_3)\to\OS(X)\) with
\(\varphi(J_3)=X\), the following are equivalent:
\begin{enumerate}
\item \(\varphi\) is completely positive;
\item \(\varphi\) is contractive;
\item \(\norm{\alpha X+\beta X^*}\le
\sqrt{|\alpha|^2+|\beta|^2}\) for all \(\alpha,\beta\in\C\).
\end{enumerate}

\begin{figure}[ht]
\centering
\includegraphics[width=0.98\linewidth]{figures/open_problem_crop.png}
\caption{Conjecture 6.10 and its extension formulation in
arXiv:2509.18027, p.~21.}
\end{figure}

The conjecture is false. In fact, condition (iii) need not imply even
3-positivity.

\section*{Proof Intuition}

The source paper already contains the obstruction, but it must be viewed
through matrix-convex duality. Its Example 4.11 gives a matrix \(B_0\) for
which \(B_0B_0^*+B_0^*B_0\le I\), yet \(B_0\) cannot be compressed from any
matrix satisfying equality. The equality matrices
\[
  G_{a,b}=\begin{bmatrix}0&a\\b&0\end{bmatrix},
  \qquad |a|^2+|b|^2=1,
\]
are exactly the elementary boundary tests used by the scalar norm condition.

Take the closed matrix-convex hull of these \(G_{a,b}\). The matrix \(B_0\)
does not belong to it, since membership would itself give an equality
dilation. The Effros--Winkler matrix Hahn--Banach theorem therefore produces
a monic pencil which is positive on every \(G_{a,b}\) but not at \(B_0\).
Writing the two self-adjoint coefficients of this pencil as one operator
\(X\), positivity on all \(G_{a,b}\) becomes exactly
\(\norm{\alpha X+\beta X^*}\le\sqrt{|\alpha|^2+|\beta|^2}\).
Failure at \(B_0\), however, says that the corresponding map on
\(\OS(J_3)\) is not 3-positive.

\section*{The Counterexample}

\begin{theorem}\label{thm:counterexample}
There exists \(X\in\M_3\) such that
\begin{equation}\label{eq:norm-condition}
 \norm{\alpha X+\beta X^*}
 \le \sqrt{|\alpha|^2+|\beta|^2}
 \qquad(\alpha,\beta\in\C),
\end{equation}
but the unital self-adjoint preserving map
\[
 \varphi:\OS(J_3)\longrightarrow\OS(X),\qquad \varphi(J_3)=X,
\]
is not 3-positive. Consequently \(\varphi\) is contractive but is not
completely positive, and Conjecture 6.10 of \cite{DDM} is false.
\end{theorem}

\begin{lemma}\label{lem:no-dilation}
Let
\[
B_0=
\begin{bmatrix}
0&0&1/\sqrt3\\
\sqrt{2/3}&0&0\\
0&1/\sqrt3&0
\end{bmatrix}.
\]
Then
\[
B_0B_0^*+B_0^*B_0=\operatorname{diag}(1,1,2/3),
\]
and there is no Hilbert space \(\mathcal L\supset\C^3\) and
\(C\in\mathcal B(\mathcal L)\) such that
\[
 P_{\C^3}C|_{\C^3}=B_0,
 \qquad CC^*+C^*C=I_{\mathcal L}.
\]
\end{lemma}

\begin{proof}
The displayed identity follows by multiplication. Suppose such a \(C\)
exists and write, relative to \(\mathcal L=\C^3\oplus\mathcal L_0\),
\[
C=\begin{bmatrix}B_0&A\\D&E\end{bmatrix}.
\]
The upper-left corner of \(CC^*+C^*C=I\) gives
\begin{equation}\label{eq:defect}
 AA^*+D^*D
 =I_3-(B_0B_0^*+B_0^*B_0)
 =\operatorname{diag}(0,0,1/3).
\end{equation}
Hence \(A=e_3a\) and \(D=de_3^*\) for some row operator \(a\) and
column vector \(d\) over \(\mathcal L_0\). The upper-right corner is
\begin{equation}\label{eq:offdiag}
 B_0D^*+AE^*+B_0^*A+D^*E=0.
\end{equation}
But
\[
 B_0e_3=\frac1{\sqrt3}e_1,
 \qquad B_0^*e_3=\frac1{\sqrt3}e_2.
\]
The \(e_1\)-row of \eqref{eq:offdiag} is \(d^*/\sqrt3=0\), and its
\(e_2\)-row is \(a/\sqrt3=0\); the two remaining terms have range in
\(\C e_3\). Thus \(a=d=0\), contradicting \eqref{eq:defect}.
\end{proof}

\begin{lemma}\label{lem:compact}
Let
\[
 \G=\left\{G_{a,b}=\begin{bmatrix}0&a\\b&0\end{bmatrix}:
 |a|^2+|b|^2=1\right\}\subset\M_2,
\]
and let \(\K=\mconv(\G)\), viewed as a matrix-convex set of pairs
\((\Rea B,\Ima B)\). Then every level \(\K(n)\) is compact, and
\(B_0\notin\K(3)\).
\end{lemma}

\begin{proof}
At level \(n\), the matrix-convex hull consists of
\begin{equation}\label{eq:mconv-rep}
 B=\sum_{j=1}^N V_j^*G_jV_j,
 \qquad G_j\in\G,\quad V_j:\C^n\to\C^2,\quad
 \sum_{j=1}^N V_j^*V_j=I_n.
\end{equation}
This set is compact. For completeness, normalize every nonzero \(V\) by
\(\operatorname{tr}(V^*V)=1\) and consider the compact set of pairs
\[
 (V^*GV,V^*V)\in \M_n\oplus(\M_n)_{\rm sa}.
\]
Equation \eqref{eq:mconv-rep} says that \((B,I_n)\) belongs to the
finite-dimensional conic hull of this compact set. Conic Carath\'eodory
reduces every representation to a uniformly bounded number of terms, while
the identity constraint fixes the total coefficient sum to
\(\operatorname{tr}I_n=n\). The resulting parameter space is compact, so
its image \(\K(n)\) is compact.

If \(B_0\) had a representation \eqref{eq:mconv-rep}, set
\(C=\bigoplus_jG_j\) and \(Vh=(V_1h,\ldots,V_Nh)\). Then \(V\) is an
isometry, \(V^*CV=B_0\), and
\[
 CC^*+C^*C=\bigoplus_j(G_jG_j^*+G_j^*G_j)=I.
\]
This contradicts Lemma \ref{lem:no-dilation}.
\end{proof}

\begin{proof}[Proof of Theorem \ref{thm:counterexample}]
The matrix-convex set \(\K\) is closed and contains zero: compressing any
\(G_{a,b}\) to its first coordinate gives zero at level one, and direct sums
give zero at every level. Apply the
Effros--Winkler matrix Hahn--Banach theorem \cite[Theorem 2.2]{HKM} at
level three to the self-adjoint pair
\((\Rea B_0,\Ima B_0)\notin\K(3)\). After changing signs of the
coefficients if necessary, it gives \(A_1,A_2\in(\M_3)_{\rm sa}\) such that
the monic pencil
\begin{equation}\label{eq:real-pencil}
 \mathcal L(B)=I+\Rea(B)\otimes A_1+\Ima(B)\otimes A_2
\end{equation}
is positive semidefinite for every \(B\in\K\), but
\(\mathcal L(B_0)\not\ge0\).

Set \(X=(A_1-iA_2)/2\in\M_3\). For every complex matrix \(B\), direct
expansion gives
\begin{equation}\label{eq:complex-pencil}
 \mathcal L(B)=I+B\otimes X+B^*\otimes X^*.
\end{equation}
For \(G_{a,b}\in\G\), positivity of \eqref{eq:complex-pencil} reads
\[
 \begin{bmatrix}
 I_3&aX+\overline bX^*\\
 bX+\overline aX^*&I_3
 \end{bmatrix}\ge0.
\]
By the Schur complement, this is equivalent to
\begin{equation}\label{eq:boundary-norm}
 \norm{aX+\overline bX^*}\le1
 \qquad(|a|^2+|b|^2=1).
\end{equation}
Given arbitrary \(\alpha,\beta\in\C\), normalize
\(a=\alpha/r\), \(b=\overline\beta/r\), where
\(r=\sqrt{|\alpha|^2+|\beta|^2}\). Equation
\eqref{eq:boundary-norm} gives \eqref{eq:norm-condition}; the zero case is
immediate.

It remains to show that the resulting map is not completely positive.
Lemma 5.2 of \cite{DDM} states that
\[
 I_n\otimes I_3+B\otimes J_3+B^*\otimes J_3^*\ge0
 \quad\Longleftrightarrow\quad BB^*+B^*B\le I_n.
\]
For \(B=B_0\), the right side holds by Lemma \ref{lem:no-dilation}.
Thus the source pencil is positive in \(\M_3(\OS(J_3))\). Applying
\(\operatorname{id}_{\M_3}\otimes\varphi\) produces
\[
 I_3\otimes I_3+B_0\otimes X+B_0^*\otimes X^*=\mathcal L(B_0),
\]
up to the harmless tensor-factor flip. This is not positive by construction.
Hence \(\varphi\) is not 3-positive and is not completely positive.

Finally, Theorem 5.1 of \cite{DDM} proves that
\eqref{eq:norm-condition} is equivalent to contractivity of \(\varphi\).
Therefore \(\varphi\) is contractive but not completely positive.
\end{proof}

\section*{Verification Notes}

The script \texttt{code/check\_finite\_identities.py} verifies the defect of
\(B_0\), equality for sample boundary generators, and the two saturated
directions \(B_0e_3=e_1/\sqrt3\) and \(B_0^*e_3=e_2/\sqrt3\). It is a
sanity check only. The existence of \(X\) is supplied by the cited matrix
Hahn--Banach theorem, not by numerical optimization.

The decisive external statement was checked in
\cite[Theorem 2.2]{HKM}: if a closed matrix-convex set containing zero omits
a point at level \(m\), a monic linear pencil of size \(m\) separates that
point from the entire set.

\section*{Scope and Novelty Check}

This existential finite-dimensional counterexample disproves the implication
(iii) \(\Rightarrow\) (i), and hence the conjectured equivalence. It does not
give numerical entries for \(X\), nor characterize the full matricial range
of \(J_3\).

The run indexes were searched for arXiv:2509.18027, the title, ``\(J_3\)'',
``matricial range'', and ``unital contractive map''. A bounded web search on
2026-07-22 used the exact phrases ``Conjecture 6.10'', ``Is every unital
contractive map from \(\OS(J_3)\)'', the source title plus ``counterexample'',
and the arXiv id plus ``Effros Winkler''. It found the source preprint and
summaries but no later proof, disproof, or discussion of this argument.

\section*{Human Review Recommendation}

Recommended verdict: likely valid counterexample. The two main review points
are Lemma \ref{lem:compact}, which ensures separation applies, and the
conversion \eqref{eq:real-pencil}--\eqref{eq:complex-pencil}. The dilation
obstruction is proved independently rather than taken from the source's brief
``analogous argument'' for arbitrary dilation size.

\begin{thebibliography}{9}
\bibitem{DDM}
Pankaj Dey, Atanu Dhang, and Mithun Mukherjee,
\emph{Matricial ranges, dilations, and unital contractive maps},
arXiv:2509.18027, 2025.

\bibitem{EW}
Edward G. Effros and S\o ren Winkler,
\emph{Matrix convexity: operator analogues of the bipolar and Hahn--Banach
theorems}, Journal of Functional Analysis \textbf{144} (1997), 117--152,
doi:10.1006/jfan.1996.2958.

\bibitem{HKM}
J. William Helton, Igor Klep, and Scott McCullough,
\emph{The tracial Hahn--Banach theorem, polar duals, matrix convex sets, and
projections of free spectrahedra}, Analysis \& PDE \textbf{11} (2018),
2343--2400; arXiv:1407.8198. Theorem 2.2 restates the Effros--Winkler
separation theorem in the form used here.
\end{thebibliography}

\end{document}
